\documentclass[border=4pt]{standalone}
\usepackage{mathptmx}
\usepackage[T1]{fontenc}
\usepackage{graphicx}
\usepackage{tikz}
\usetikzlibrary{shapes.geometric,shapes.misc,arrows.meta,positioning,calc,fit,shadows.blur}
\usepackage{xcolor}

\definecolor{diagLaw}{HTML}{1D6FE0}
\definecolor{diagPlan}{HTML}{E63946}
\definecolor{diagEnv}{HTML}{12A150}
\definecolor{diagLatent}{HTML}{7C3AED}
\definecolor{diagProbe}{HTML}{E8890C}
\definecolor{gapblue}{HTML}{00B4D8}
\definecolor{gapred}{HTML}{C2185B}
\definecolor{inkmute}{HTML}{6B7280}

\begin{document}
\begin{tikzpicture}[
    >=Stealth,
    font=\sffamily\small,
    % Flat fill + a heavy tinted border carries node identity; the blur shadow is barely there
    % (opacity 11) so a card lifts off the page without reading as a 2005 gradient. Rounded
    % rectangles also edge-align, which is what lets the three rows line up -- ellipses cannot.
    card/.style={draw, line width=1.5pt, rounded corners=3pt, align=center, fill=white,
                 minimum height=10mm, inner xsep=2.6mm, font=\sffamily\large\bfseries,
                 blur shadow={shadow blur steps=8, shadow xshift=0.5pt, shadow yshift=-0.8pt,
                              shadow opacity=11, shadow blur radius=1.3pt}},
    pic/.style={card, inner sep=1pt, minimum height=0pt},       % image cards carry no text padding
    lat/.style ={card, draw=diagLatent!70!black, text=diagLatent!25!black},
    prb/.style ={card, draw=diagProbe!75!black,  text=diagProbe!30!black},
    law/.style ={card, draw=diagLaw!70!black,    text=diagLaw!25!black},
    pln/.style ={card, draw=diagPlan!70!black,   text=diagPlan!25!black},
    % Ordinary dataflow is deliberately recessive: thin and grey. The gaps must win the page.
    flow/.style={-{Stealth[length=5.5pt]}, line width=0.9pt, draw=inkmute!60},
    % THE MOTIF: the two isomorphism gaps are the conceptual punchline, so they are the heaviest
    % marks here and the only two horizontal arrows in the middle of the figure -- mirrored, one
    % pointing right into the theory, one pointing left back out into control.
    gap/.style={line width=2.4pt, dash pattern=on 5.5pt off 3pt, -{Stealth[length=7pt]}},
    % Solid-colour capsule, white text, sitting just ABOVE the line it names so the dashes show.
    caps/.style={fill=#1, text=white, rounded corners=4pt, inner xsep=4.5pt, inner ysep=2.2pt,
                 font=\sffamily\normalsize\bfseries},
    rail/.style={font=\sffamily\normalsize\bfseries, text=inkmute!75, rotate=90, anchor=center},
    plab/.style={font=\sffamily\normalsize\bfseries, align=center},
]

%======================= column 1: the environment, spanning law + perception =======================%
\node[pic, draw=diagEnv!70!black] (world) at (2.6,-0.15)
      {\includegraphics[width=38mm]{inset_world_p.png}};
\node[plab, text=diagEnv!30!black, anchor=south] at (2.6,2.92) {Isaac Lab pusher};

%======================= LAW row: the source text, and the CLASSICAL isomorphism =======================%
\node[law] (statute) at (8.0,2.4)  {Legal source text};
\node[law] (rules)   at (16.0,2.4) {DDL rule base};
\draw[gap, draw=inkmute!60, line width=1.7pt, dash pattern=on 3.5pt off 2.4pt]
      (statute) -- node[caps=inkmute!70, above=0.4mm] {isomorphism} (rules);

%======================= PERCEPTION row: world -> atoms  (the GROUNDING gap) =======================%
\node[lat] (wm)     at (7.0,0.2)  {World model $\psi$};
\node[prb] (probes) at (10.5,0.2)  {Probes $\phi$};
\node[law] (atoms)  at (16.0,0.2) {Atoms};
\draw[flow] (world.east) -- (wm.west);
\draw[flow] (wm) -- (probes);
\draw[gap, draw=gapblue] (probes) -- node[caps=gapblue, above=0.4mm] {Grounding gap} (atoms);

%======================= CONTROL row: verdict -> constraint  (the ONTOLOGICAL gap) =======================%
\node[law] (verdict) at (16.0,-2.3) {Verdict};
\node[pln] (cons)    at (10.3,-2.3)  {Planning constraint};
\node[pic, draw=diagPlan!70!black] (planner) at (6.9,-2.3)
      {\includegraphics[width=17mm]{inset_planner.png}};
\node[plab, text=diagPlan!25!black, anchor=north] at (6.9,-3.25) {RRT tree};
\draw[flow] (rules) -- (atoms);
\draw[flow] (atoms) -- (verdict);
\draw[gap, draw=gapred] (verdict) -- node[caps=gapred, above=0.4mm] {Ontological gap} (cons);
\draw[flow] (cons.west) -- (planner.east);

%======================= the closed loop, made explicit =======================%
\draw[flow, line width=1.0pt] (planner.west) -- (world.east|-planner.west);
\node[font=\sffamily\normalsize, text=inkmute, anchor=south] at (5.35,-2.22) {action};

%======================= the left rail: the three layers of the stack =======================%
\node[rail] at (0.35,2.4)  {\textsc{law}};
\node[rail] at (0.35,0.2)  {\textsc{perception}};
\node[rail] at (0.35,-2.3) {\textsc{control}};

\end{tikzpicture}
\end{document}
